\chapter{Архитектура и среда разработки}
\label{ch:architecture}

\section{Общая архитектура системы}

Видеохостинг состоит из четырёх взаимодействующих частей: обратного
прокси, клиентского приложения, серверного приложения и хранилища
данных. Запрос пользователя проходит следующий путь:

\begin{lstlisting}[breaklines=false]
            [ Browser ]
                 |  HTTPS
                 v
            [ Caddy ]   <-- reverse proxy + TLS
                 |
                 v
            [ nginx ]   <-- SPA, /media, /static
                 |  /api/*
                 v
        [ Django + Gunicorn ]
                 |
                 v
          [ PostgreSQL ]
\end{lstlisting}

\textbf{Caddy} принимает внешние HTTPS-запросы и автоматически выпускает
TLS-сертификат Let's~Encrypt. \textbf{nginx} отдаёт собранное
React-приложение и статические файлы (видео, превью), а обращения к
\texttt{/api/} проксирует на серверное приложение. \textbf{Django} с
WSGI-сервером \textbf{Gunicorn} обслуживает REST-API, обращаясь к
\textbf{PostgreSQL} за данными.

Серверное приложение разделено на три Django-приложения по принципу
единственной ответственности:
\begin{itemize}
  \item \texttt{accounts} — пользователи, аутентификация, гостевые
        аккаунты, выдача JWT;
  \item \texttt{videos} — видео, загрузка, потоковая отдача;
  \item \texttt{common} — общий код: базовые модели, разрешения,
        постраничная навигация, отдача файлов с поддержкой Range.
\end{itemize}

\section{Модель данных}

База данных описана классами-моделями Django. Основные сущности:

\begin{itemize}
  \item \textbf{User} — пользователь. Идентификация по имени
        пользователя, вход по паролю. Хранит необязательный e-mail,
        видимое имя, признак гостевого аккаунта и отметку времени
        последней активности.
  \item \textbf{Video} — видеоролик: владелец, название, описание,
        файл, превью, длительность, размер, число просмотров, признак
        публичности. Первичный ключ — \texttt{UUID}.
\end{itemize}

Общие поля (отметки создания и изменения, UUID-ключ) вынесены в
абстрактные базовые модели приложения \texttt{common} и наследуются
остальными моделями — это исключает дублирование кода (принцип DRY).

\section{REST-API}

REST-API смонтирован под префиксом \texttt{/api/} и сгруппирован по
областям:

\begin{itemize}
  \item \textbf{Авторизация} (\texttt{/api/auth/}): регистрация, вход по
        паролю, гостевой вход, обновление токена, чтение и правка
        профиля, выход со всех устройств.
  \item \textbf{Видео} (\texttt{/api/videos/}): лента, загрузка,
        карточка ролика, потоковая отдача файла, список своих видео,
        правка и удаление.
\end{itemize}

Полное описание методов с примерами запросов оформлено в репозитории в
файле \texttt{API.md}.

\section{Окружение разработки и развёртывания}

Конфигурация приложения управляется переменными окружения. Это
позволяет одному и тому же коду работать в двух режимах:
\begin{itemize}
  \item \textbf{локальная разработка} — без внешних сервисов, база
        данных \texttt{SQLite};
  \item \textbf{продакшен} — \texttt{PostgreSQL}; значения задаются в
        файле \texttt{.env}, который не попадает в репозиторий.
\end{itemize}

Развёртывание автоматизировано. Образы серверной и клиентской частей
описаны в \texttt{Dockerfile}; файл \texttt{docker-compose.yml}
поднимает весь стек одной командой. Скрипт \texttt{deploy.sh}
разворачивает платформу на сервере с уже настроенным Caddy: он
синхронизирует исходники, добавляет сервисы в общий
\texttt{docker-compose.yml}, дописывает поддомен в конфигурацию Caddy и
поднимает контейнеры. Скрипт идемпотентен, повторный запуск не ломает
уже развёрнутую систему.

\endinput
